\section{Multi-Morphology Text-Associated Motion Dataset}\label{sec:dataset}

We construct a derived corpus in which a source clip identity connects a motion
description, 128 generated motion variants, and their matching simulation
assets. A \emph{clip} denotes a source segment, including a mirrored segment
when present; a \emph{variant} denotes that clip on one body configuration.
These units must not be conflated when reporting dataset size or evaluation
coverage. Figure~\ref{fig:data_pipeline} shows the construction paths.

\begin{figure}[t]
\centering
\fbox{\begin{minipage}{0.94\linewidth}
\small
\begin{tabularx}{\linewidth}{@{}p{0.12\linewidth}Y@{}}
\textbf{Motion} &
AMASS $\longrightarrow$ HumanML3D preprocessing and clip extraction
$\longrightarrow$ HUMOS conditioned on $(g,\boldsymbol{\beta})$
$\longrightarrow$ MotionLib conversion and reference refinement
$\longrightarrow$ per-clip motion variants. \\[0.7em]
\textbf{Body} &
SMPL with the same $(g,\boldsymbol{\beta})$
$\longrightarrow$ SMPLSim
$\longrightarrow$ matching skeleton, collision geometry, mass/inertia,
and MJCF asset identity. \\[0.7em]
\textbf{Text} &
HumanML3D source descriptions
$\longrightarrow$ clip-ID lookup associated with the generated variants
(metadata only; not a policy input).
\end{tabularx}

\smallskip
The body assets are used for motion conversion, shape-specific forward
kinematics, ground-clearance correction, and simulation. Refined reference
motion and matching simulated state feed one shared policy.
\end{minipage}}
\caption{Dataset and asset construction. The body parameters are shared across
the motion and body paths; text follows a separate identity-association path.}\label{fig:data_pipeline}
\end{figure}

\begin{figure}[t]
\centering
\begin{tikzpicture}[>=latex, node distance=0.7cm]
\node[draw, rounded corners, fill=blue!8, minimum width=2.4cm, minimum height=0.8cm]
  (source) {20,951 clip identities};
\node[draw, rounded corners, fill=green!10, right=of source, minimum width=2.4cm, minimum height=0.8cm]
  (variants) {$\times$ 128 bodies};
\node[draw, rounded corners, fill=orange!12, right=of variants, minimum width=2.4cm, minimum height=0.8cm]
  (total) {2,681,728 variants};
\draw[->, thick] (source) -- (variants);
\draw[->, thick] (variants) -- (total);
\node[draw, rounded corners, below=0.75cm of variants, minimum width=2.1cm, minimum height=0.65cm]
  (train) {18,855 train};
\node[draw, rounded corners, right=0.35cm of train, minimum width=1.8cm, minimum height=0.65cm]
  (val) {1,048 val};
\node[draw, rounded corners, right=0.35cm of val, minimum width=1.8cm, minimum height=0.65cm]
  (test) {1,048 test};
\draw[->, thick] (source.south) |- (train.north);
\draw[->, thick] (source.south) |- (val.north);
\draw[->, thick] (source.south) |- (test.north);
\end{tikzpicture}
\caption{Scale and split structure of the derived corpus. Splitting is by clip
identity, so all 128 body variants of a clip remain in the same partition.}
\label{fig:dataset_scale}
\end{figure}

\subsection{Source motion and text provenance}\label{sec:source_data}

We use the AMASS-derived motion portion of HumanML3D~\cite{amass,humanml3d}, with the 18 source collections listed in
Appendix~\ref{app:preprocessing}. Source SMPL+H sequences are processed at
20\,Hz, including the Z-up to Y-up coordinate conversion used by the
HumanML3D/HUMOS preprocessing pipeline. Clip extraction, temporal cropping, and
mirroring preserve the HumanML3D identities. The preparation also includes
treadmill/skating cleanup. This subset is not the entirety of AMASS.

The generation pipeline produced 22,459 candidate clip identities, counting
original and mirrored clips. A subsequent text-based filter excluded 1,508
clips identified as requiring external objects or environmental support,
leaving 20,951 identities for the flat-ground imitation task. This filter is a
task-scope heuristic: it does not prove that every retained reference is
support-free or dynamically feasible. The retained and excluded ID lists
provide the audit trail for that decision.

Source descriptions are available through a separate clip-ID lookup and the
HumanML3D annotation metadata used by HUMOS. A description is associated with
the variants of its source clip; it is not a new human annotation for each
generated body. A coarse coverage check against a single-caption-per-clip
lookup finds a caption for all 20,951 clip identities, but semantic fidelity
is not established. Of the 9,579 mirrored identities, 22\% carry a caption
byte-identical to their un-mirrored source, and 40\% contain an explicit
left/right or clockwise term that mirroring inverts; simply removing a mirror
prefix therefore does not establish a correct caption. Segment timing and
multi-annotation coverage against the full HumanML3D annotation set are not
audited here.

\subsection{Body sampling and shape-specific generation}\label{sec:body_generation}

Let a body configuration be $b=(g,\boldsymbol{\beta})$, where $g$ identifies the
female or male SMPL body-model variant and
$\boldsymbol{\beta}\in\R^{10}$ contains shape coefficients. We sample
\begin{equation}
    \beta_{ij}\sim\mathcal{U}(-3,3),\qquad
    i=1,\ldots,64,\quad j=1,\ldots,10,
    \label{eq:shape_sampling}
\end{equation}
using NumPy's \texttt{default\_rng(46)}, and cross these 64 vectors with the two
body-model variants. Thus the corpus contains 128 body configurations, not
128 independently sampled beta vectors. The actual coefficient arrays and
stable beta keys identify the bodies; the keys are not content hashes.
Uniform sampling of model coefficients is a controlled source of variation,
not a representative sample of human populations.

For each source clip, we provide the processed pose/root motion features and
the target body parameters to the pretrained HUMOS
model~\cite{humos}. The local configuration uses root-relative 6D root and
joint rotations together with translation, with shape and body-model
conditioning. The generated pose and root sequences are retained for
conversion. HUMOS is an existing identity-conditioned model; neither its
training nor a text-to-motion generator is a contribution of this work.
The checkpoint used in the documented pipeline is identified in
Appendix~\ref{app:provenance}.

In parallel, SMPL bodies~\cite{smpl} with identical parameters are converted
to MJCF assets using SMPLSim~\cite{smplsim}. The asset-generation configuration
uses mass estimation and the configured hand/foot simplifications
(\texttt{real\_weight}, \texttt{replace\_feet}, \texttt{remove\_toe}, and
\texttt{freeze\_hand}). The resulting imitation representation contains
24 rigid bodies and 69 actuated joint coordinates. Collision proxies, joint
ranges, and estimated masses/inertias belong to each asset. These physical
parameters are simulation-model estimates, not measurements of corresponding
human participants.

\subsection{Motion conversion and reference refinement}\label{sec:refinement}

The generated sequences are converted to ProtoMotions' Z-up simulation
convention and resampled from 20 to 30\,Hz. Forward kinematics (FK) uses the
matching asset rather than a single neutral skeleton. The packaged
representation includes global body positions and rotations, local rotations,
joint positions, linear/angular and joint velocities, and contact labels.
The earlier corpus already included frame-zero grounding; the newer
refinement addresses artifacts throughout each sequence.

Refinement is applied independently to each variant in the following order.
First, stored ankle/toe contacts are validated, or contacts are detected from
foot position and velocity if usable labels are unavailable. Short gaps are
closed and very short contact runs are removed. These labels define candidate
stance intervals; they are not measured ground-reaction forces.

Second, joint rotations are temporally filtered using sign-continuous
quaternion representations, normalization, and a bounded geodesic correction.
Contact-transition frames divide the sequence into segments so filtering does
not average across a change of stance. A transition taper and a limited
smoothing blend further restrict changes near those boundaries.
Post-smoothing FK then recomputes the positions used by subsequent corrections.

Third, a conservative root-translation correction reduces horizontal stance
sliding. For each contact interval, an anchor is formed from the foot's median
horizontal position. A least-squares correction seeks a common root XY shift
that brings the active feet closer to their anchors; the shift is temporally
regularized and capped. Multiple contacts can disagree, so this root-only
operation does not guarantee exact foot locking and does not solve for new
joint angles.

Fourth, we evaluate the lowest point of the matching collision geometry at
every frame. With floor height zero and margin $\epsilon_z$, the required lift is
\begin{equation}
    \ell_t^{\mathrm{req}}
      = \max\!\left(0,\epsilon_z-\min_{x\in\mathcal{C}_b(t)}x_z\right),
    \label{eq:clearance}
\end{equation}
where $\mathcal{C}_b(t)$ denotes the modeled collision geometry for body $b$.
A smoothed lift envelope is constrained to remain above this requirement.
This prevents penetration of the checked geometry in the reference sequence;
it does not establish balanced support, realistic impacts, or freedom from
self-collision.

Finally, FK and velocities are recomputed from the corrected root and
rotations, independently within each motion. This avoids inconsistent
position/velocity tensors and artificial derivatives across clip boundaries.
The default thresholds, filter windows, and correction bounds are given in
Table~\ref{tab:refinement_parameters}.

The refinement objective is modest: reduce visible and measurable kinematic
artifacts while preserving the generated motion. Before/after assessment must
therefore include both artifact measures (sliding, penetration, and jerk) and
fidelity measures (rotation/root changes and retained cross-shape variation).
Improved kinematics alone does not demonstrate improved policy learning or
dynamic feasibility.

\subsection{Packaging, scale, and deterministic splits}\label{sec:splits}

Each per-clip file contains the 128 variants, their frame counts and sampling
rates, and the clip/asset identities needed to recover the corresponding body
parameters. Captions remain a separately associated metadata resource; they
are not established as a field of every training tensor file. This packaging
allows the training system to stream clip files without downloading the
complete corpus at startup.

\begin{table}[t]
\centering
\caption{Current corpus and split sizes. Variants are generated combinations,
not independent captures. Every clip has 128 body configurations.}\label{tab:dataset_counts}
\begin{tabular}{lrr}
\toprule
Partition & Clip identities & Motion variants \\
\midrule
Training & 18,855 & 2,413,440 \\
Validation & 1,048 & 134,144 \\
Test & 1,048 & 134,144 \\
\midrule
Total & 20,951 & 2,681,728 \\
\bottomrule
\end{tabular}
\end{table}

We use the versioned \texttt{hhi\_stage2\_v1} split with seed 42 and approximate
90/5/5 proportions (Table~\ref{tab:dataset_counts}). All variants of a clip
stay together, and original/mirrored partners are assigned as a group.
The 150 development clips are reserved for training, together with their
mirror partners. Training accepts explicit training and validation manifests;
test-manifest membership is not loaded by the training workflow. Split hashes
and version information accompany the run and checkpoints.

These splits test held-out clip identities on the trained body set. They have
not been demonstrated to be disjoint by subject or original capture recording:
different segments can still share a source recording. This distinction
limits claims about broader motion generalization.

\subsection{Current resource scope}

The current resource combines generated references, matching assets, inherited
text associations, and reproducible splits. It is not a fully audited
text-and-motion benchmark release. Caption completeness and alignment,
body-specific language, licensing and redistribution conditions, and
shape-conditioned semantic fidelity require further work. Those extensions
can be developed independently of the reference-conditioned controller
evaluated here.
